
\documentclass[12pt,a4paper]{article}

% =========================================================
% Кодировка и язык
% =========================================================
\usepackage[T2A]{fontenc}
\usepackage[utf8]{inputenc}
\usepackage[russian,english]{babel}

% =========================================================
% Геометрия и интервалы
% =========================================================
\usepackage{geometry}
\geometry{left=2.5cm,right=2cm,top=2cm,bottom=2cm}
\usepackage{setspace}
\onehalfspacing

% =========================================================
% Математика, таблицы, графика
% =========================================================
\usepackage{amsmath,amssymb}
\usepackage{booktabs}
\usepackage{longtable}
\usepackage{array}
\usepackage{multirow}
\usepackage{graphicx}
\usepackage{xcolor}
\usepackage{float}
\usepackage{caption}
\usepackage{enumitem}

% =========================================================
% Листинги
% =========================================================
\usepackage{listings}
\lstdefinestyle{code}{
    basicstyle=\ttfamily\footnotesize,
    commentstyle=\color{gray!70!black}\itshape,
    numbers=left,
    numberstyle=\tiny\color{gray},
    numbersep=8pt,
    frame=single,
    rulecolor=\color{gray!50},
    breaklines=true,
    showstringspaces=false,
    tabsize=4,
    columns=fullflexible,
    keepspaces=true,
    backgroundcolor=\color{gray!5},
}
\lstset{style=code}

% =========================================================
% Ссылки
% =========================================================
\usepackage{hyperref}
\hypersetup{
    colorlinks=true,
    linkcolor=blue,
    urlcolor=blue
}

% =========================================================
% Абзацы
% =========================================================
\setlength{\parindent}{1.25cm}
\setlength{\parskip}{0pt}

% =========================================================
% Документ
% =========================================================
\begin{document}

% =========================================================
% Титульный лист
% =========================================================
\begin{titlepage}
\begin{center}

САНКТ-ПЕТЕРБУРГСКИЙ ГОСУДАРСТВЕННЫЙ УНИВЕРСИТЕТ

\vspace{1cm}

Факультет математики и механики

\vfill

{\Large \textbf{ОТЧЁТ ПО УЧЕБНОЙ ПРАКТИКЕ}}

\vspace{0.5cm}

{\large \textbf{«Сравнение и анализ оптимизаций компиляторов»}}

\vfill

\begin{flushright}
Выполнил: студент группы \underline{\hspace{3cm}}\\
\underline{\hspace{6cm}}\\
ФИО

\vspace{0.7cm}

Научный руководитель:\\
\underline{\hspace{6cm}}\\
ФИО
\end{flushright}

\vfill

Санкт-Петербург\\
2026

\end{center}
\end{titlepage}

\tableofcontents
\newpage

% =========================================================
\section{Введение}
% =========================================================

Современные компиляторы языка C выполняют большое количество преобразований исходной программы, направленных на повышение скорости выполнения, уменьшение объёма машинного кода и более эффективное использование ресурсов процессора.

Один и тот же исходный код при использовании различных компиляторов и разных уровней оптимизации может быть преобразован в существенно отличающийся машинный код. Особенно важен выбор стратегии оптимизации при разработке системного и встраиваемого программного обеспечения, где одновременно предъявляются требования к быстродействию и размеру исполняемого кода.

Цель работы --- исследовать и сравнить методы оптимизации, применяемые компиляторами языка C GCC, Clang и Tiny C Compiler, а также определить влияние оптимизации на ассемблерный код, размер программы и время её выполнения.

Для достижения цели необходимо решить следующие задачи:
\begin{enumerate}
    \item выбрать основной компилятор;
    \item определить сценарий применения языка C;
    \item провести обзор компиляторов, подходящих для выбранного сценария;
    \item выбрать дополнительные компиляторы для сравнения;
    \item зафиксировать версии компиляторов и характеристики системы;
    \item исследовать программу \texttt{optbench.c};
    \item сравнить размер и время работы программ без оптимизации, с оптимизацией по скорости и по размеру;
    \item сравнить ассемблерный код до и после оптимизации;
    \item проанализировать удалённые и заменённые инструкции;
    \item сравнить стратегии оптимизации различных компиляторов;
    \item разработать собственные тесты для выбранного сценария;
    \item провести аналогичный эксперимент на собственных тестах;
    \item сформулировать выводы.
\end{enumerate}

% =========================================================
\section{Выбор сценария применения языка C}
% =========================================================

В качестве рассматриваемого сценария выбрана разработка встраиваемого программного обеспечения для промышленных контроллеров и других ресурсоограниченных систем.

Во встраиваемых системах программа обычно работает в условиях ограниченного количества оперативной и постоянной памяти. Для промышленных контроллеров также важна предсказуемость времени выполнения отдельных участков программы.

Язык C широко применяется в подобных задачах благодаря следующим особенностям:
\begin{itemize}
    \item возможности прямой работы с памятью;
    \item небольшим накладным расходам во время выполнения;
    \item поддержке работы с аппаратными регистрами;
    \item наличию битовых операций;
    \item отсутствию обязательной среды выполнения высокого уровня;
    \item доступности компиляторов для большого количества архитектур;
    \item возможности получения компактного машинного кода.
\end{itemize}

Для выбранного сценария особенно важны:
\begin{itemize}
    \item качество оптимизации циклов;
    \item оптимизация арифметических операций;
    \item эффективное использование регистров;
    \item корректная работа с \texttt{volatile};
    \item уменьшение объёма исполняемого кода;
    \item оптимизация битовых операций;
    \item устранение лишних обращений к памяти.
\end{itemize}

Основные эксперименты выполнялись на архитектуре x86-64. Дополнительные тесты при этом моделируют характерные для embedded-систем операции: работу с регистрами, кольцевыми буферами, CRC и целочисленной фильтрацией.

% =========================================================
\section{Выбор компиляторов}
% =========================================================

\subsection{GCC}

В качестве основного компилятора выбран \textbf{GCC (GNU Compiler Collection)}.

GCC является одним из наиболее распространённых свободных компиляторов языка C, поддерживает большое количество аппаратных архитектур и имеет развитую систему оптимизационных проходов.

В работе рассматриваются стандартные режимы:
\begin{itemize}
    \item \texttt{-O0} --- практически отключённая оптимизация;
    \item \texttt{-O2} --- стандартная оптимизация производительности;
    \item \texttt{-Os} --- оптимизация с дополнительным приоритетом размера кода.
\end{itemize}

\subsection{Clang}

В качестве второго компилятора выбран \textbf{Clang 18.1.3}, использующий инфраструктуру LLVM.

Clang предоставляет современные средства оптимизации, включая оптимизацию циклов, устранение общих подвыражений, векторизацию и межпроцедурные преобразования.

В эксперименте использовались уровни \texttt{-O0}, \texttt{-O2} и \texttt{-Os}.

\subsection{Tiny C Compiler}

Третьим компилятором выбран \textbf{Tiny C Compiler 0.9.27 (TCC)}.

TCC ориентирован прежде всего на высокую скорость компиляции и имеет значительно более простую систему оптимизации по сравнению с GCC и Clang. В работе он используется как контрольный вариант, позволяющий сравнить код, полученный без развитой системы глобальных оптимизационных проходов.

% =========================================================
\section{Обзор источников}
% =========================================================

Основой методики эксперимента является статья Ф.~Н.~Хислея «Генерация высококачественного кода для программ, написанных на Си».

В статье рассматривается программа \texttt{optbench.c}, состоящая из отдельных фрагментов, предназначенных для проверки различных методов оптимизации. Программа намеренно не является образцом хорошего прикладного программирования: её назначение --- создать ситуации, в которых оптимизатор может удалить или преобразовать исходные конструкции.

В частности, проверяются:
\begin{itemize}
    \item свёртка констант;
    \item распространение констант;
    \item распространение копий;
    \item арифметические упрощения;
    \item устранение общих подвыражений;
    \item устранение лишних присваиваний;
    \item снижение мощности операций;
    \item оптимизация переменных индукции;
    \item вынесение инвариантного кода из цикла;
    \item удаление ненужных циклов;
    \item устранение недостижимого кода;
    \item разворачивание циклов;
    \item слияние циклов;
    \item сжатие цепочек переходов;
    \item использование регистров;
    \item машинно-зависимые оптимизации.
\end{itemize}

Дополнительно использовались официальная документация GCC, Clang и TCC, а также материалы по инструментам GNU Binutils.

% =========================================================
\section{Характеристики системы и версии компиляторов}
% =========================================================

\begin{table}[H]
\centering
\caption{Параметры системы}
\begin{tabular}{ll}
\toprule
Параметр & Значение \\
\midrule
ОС & Linux Mint 22.2 (Zara), база Ubuntu 24.04 \\
Архитектура & x86\_64 \\
GCC & 13.3.0 \\
Clang & 18.1.3 \\
TCC & 0.9.27 \\
\bottomrule
\end{tabular}
\end{table}

\begin{table}[H]
\centering
\caption{Версии исследуемых компиляторов}
\begin{tabular}{lll}
\toprule
Компилятор & Версия & Команда \\
\midrule
GCC & 13.3.0 & \texttt{gcc} \\
Clang & 18.1.3 & \texttt{clang} \\
Tiny C Compiler & 0.9.27 & \texttt{tcc} \\
\bottomrule
\end{tabular}
\end{table}

% =========================================================
\section{Методика эксперимента}
% =========================================================

\subsection{Сборка GCC}

\begin{lstlisting}[language=bash,caption={Сборка GCC}]
gcc -O0 -static -DNO_ZERO_DIVIDE=1 optbench.c -o gcc_O0
gcc -O2 -static -DNO_ZERO_DIVIDE=1 optbench.c -o gcc_O2
gcc -Os -static -DNO_ZERO_DIVIDE=1 optbench.c -o gcc_Os
gcc -O0 -DNO_ZERO_DIVIDE=1 optbench.c -o gcc_O0_dyn
\end{lstlisting}

\subsection{Сборка Clang}

\begin{lstlisting}[language=bash,caption={Сборка Clang}]
clang -O0 -static -DNO_ZERO_DIVIDE=1 optbench.c -o clang_O0
clang -O2 -static -DNO_ZERO_DIVIDE=1 optbench.c -o clang_O2
clang -Os -static -DNO_ZERO_DIVIDE=1 optbench.c -o clang_Os
clang -O0 -DNO_ZERO_DIVIDE=1 optbench.c -o clang_O0_dyn
\end{lstlisting}

\subsection{Сборка TCC}

\begin{lstlisting}[language=bash,caption={Сборка TCC}]
tcc -DNO_ZERO_DIVIDE=1 optbench.c -o tcc_O0
\end{lstlisting}

\subsection{Генерация ассемблерного кода}

Для GCC и Clang использовались параметры \texttt{-S -masm=intel}:

\begin{lstlisting}[language=bash]
gcc -O0 -S -masm=intel -fverbose-asm optbench.c -o gcc_O0.s
gcc -O2 -S -masm=intel -fverbose-asm optbench.c -o gcc_O2.s

clang -O0 -S -masm=intel -fverbose-asm optbench.c -o clang_O0.s
clang -O2 -S -masm=intel -fverbose-asm optbench.c -o clang_O2.s
\end{lstlisting}

Для TCC выполнялась компиляция в объектный файл и последующее дизассемблирование:

\begin{lstlisting}[language=bash]
tcc -DNO_ZERO_DIVIDE=1 -c optbench.c -o tcc_O0.o
objdump -d -M intel tcc_O0.o > tcc_O0.s
\end{lstlisting}

\subsection{Замеры}

Размер секции машинного кода определялся утилитой \texttt{size}. Время выполнения измерялось с помощью:

\begin{lstlisting}[language=bash]
/usr/bin/time -f "%e"
\end{lstlisting}

Каждая программа запускалась 10 раз.

Следует учитывать, что разрешения данного способа измерения недостаточно для очень коротких тестов. Поэтому значение 0.00 с означает, что фактическое время оказалось ниже точности использованного измерения.

% =========================================================
\section{Результаты основного теста}
% =========================================================

\begin{table}[H]
\centering
\caption{Размер секции \texttt{.text} и время выполнения}
\begin{tabular}{llrr}
\toprule
Компилятор & Режим & \texttt{.text}, байт & Время, с \\
\midrule
GCC 13.3.0 & \texttt{-O0} static & 670857 & 0.09 \\
GCC 13.3.0 & \texttt{-O2} static & 670217 & $<0.005$ \\
GCC 13.3.0 & \texttt{-Os} static & 670073 & $<0.005$ \\
GCC 13.3.0 & \texttt{-O0} dynamic & 3754 & 0.09 \\
Clang 18.1.3 & \texttt{-O0} static & 678257 & 0.09 \\
Clang 18.1.3 & \texttt{-O2} static & 678701 & $<0.005$ \\
Clang 18.1.3 & \texttt{-Os} static & 678541 & $<0.005$ \\
Clang 18.1.3 & \texttt{-O0} dynamic & 3471 & 0.09 \\
TCC 0.9.27 & стандартный режим & 4935 & 0.07 \\
\bottomrule
\end{tabular}
\end{table}

При статической сборке GCC и Clang в бинарный файл включается значительная часть библиотечного кода, поэтому абсолютный размер таких файлов нельзя напрямую сравнивать с TCC.

Для GCC минимальный размер среди исследованных статических сборок получен при \texttt{-Os}. Для Clang режим \texttt{-O2} увеличил размер секции \texttt{.text} относительно \texttt{-O0}, что показывает: оптимизация по скорости не обязана уменьшать объём кода.

По времени неоптимизированные GCC и Clang показали около 0.09 с, а оптимизированные варианты оказались ниже порога точного измерения используемым способом.

% =========================================================
\section{Анализ ассемблерного кода без оптимизации}
% =========================================================

\begin{table}[H]
\centering
\caption{Ассемблер без оптимизации: ключевые участки}
\small
\begin{tabular}{p{3cm}p{3.5cm}p{3.5cm}p{3.5cm}}
\toprule
Участок & GCC -O0 & Clang -O0 & TCC \\
\midrule
\texttt{dead\_code} & мёртвый код в основном сохранён &
мёртвый код сохранён & код сохранён, дополнительный переход \\
\texttt{unnecessary\_loop} & полный цикл &
полный цикл & полный цикл \\
\texttt{loop\_jamming} & два отдельных цикла &
два отдельных цикла & два отдельных цикла \\
\texttt{loop\_unrolling} & цикл сохранён &
цикл сохранён & цикл сохранён \\
\texttt{jump\_compression} & цепочки переходов &
цепочки переходов & вложенные переходы \\
\texttt{i3=1+2} & результат 3 & результат 3 & результат 3 \\
\texttt{k3=1;k3=1} & две записи &
две записи & две записи \\
\bottomrule
\end{tabular}
\end{table}

Даже при \texttt{-O0} некоторые простые преобразования могут выполняться на стадии генерации машинного кода. Поэтому \texttt{-O0} не следует понимать как буквальное преобразование каждой операции C в отдельную машинную инструкцию.

% =========================================================
\section{Анализ ассемблерного кода с оптимизацией}
% =========================================================

\subsection{Удаление мёртвого кода}

Функция \texttt{dead\_code} содержит присваивание, результат которого далее не используется, а также условие \texttt{if (0)}, которое никогда не выполняется.

При \texttt{-O2} GCC фактически оставляет только:

\begin{lstlisting}
endbr64
ret
\end{lstlisting}

Clang оставляет:

\begin{lstlisting}
ret
\end{lstlisting}

Таким образом, оба компилятора полностью устраняют бесполезное тело функции.

\subsection{Удаление ненужного цикла}

В функции \texttt{unnecessary\_loop} тело цикла записывает в \texttt{k5} одно и то же значение. GCC и Clang при \texttt{-O2} удаляют сам цикл и устанавливают конечное значение счётчика напрямую.

\subsection{Удаление лишнего присваивания}

Код:

\begin{lstlisting}[language=C]
k3 = 1;
k3 = 1;
\end{lstlisting}

при оптимизации сокращается до одной записи, поскольку первое значение никогда не читается.

\subsection{Снижение мощности}

При работе с константами умножение или деление может быть заменено более простой операцией.

Например, умножение на степень двойки может быть реализовано сдвигом:

\begin{lstlisting}
sal eax, 2
\end{lstlisting}

что соответствует умножению значения на четыре.

Для операции остатка от деления на степень двойки компилятор может использовать битовую маску. Например, для беззнакового значения:

\begin{lstlisting}
x % 64
\end{lstlisting}

эквивалентно:

\begin{lstlisting}
x & 63
\end{lstlisting}

\subsection{Сжатие переходов}

Функция \texttt{jump\_compression} содержит вложенные условия и переходы. После оптимизации часть промежуточных переходов устраняется, а поток управления становится компактнее.

% =========================================================
\section{Сравнение оптимизированного кода компиляторов}
% =========================================================

\begin{table}[H]
\centering
\caption{Сравнение оптимизированного машинного кода}
\small
\begin{tabular}{p{3cm}p{4cm}p{4cm}p{3cm}}
\toprule
Участок & GCC -O2 & Clang -O2 & TCC \\
\midrule
\texttt{dead\_code} & \texttt{endbr64; ret} &
\texttt{ret} & тело в большей степени сохранено \\
\texttt{unnecessary\_loop} & цикл удалён &
цикл удалён & цикл сохранён \\
Повторное присваивание & первая запись удалена &
первая запись удалена & сохраняется \\
\texttt{loop\_jamming} & циклы не слиты &
циклы не слиты & циклы не слиты \\
\texttt{loop\_unrolling} & цикл не развёрнут &
цикл не развёрнут & цикл сохранён \\
\texttt{jump\_compression} & переходы сокращены &
переходы сокращены & структура ближе к исходной \\
Общие подвыражения & оптимизированы &
оптимизированы & ограниченно \\
\bottomrule
\end{tabular}
\end{table}

GCC и Clang используют существенно более глубокий анализ программы по сравнению с TCC. При этом конкретный машинный код может различаться даже для семантически эквивалентных преобразований.

% =========================================================
\section{Сравнение методов оптимизации}
% =========================================================

\begin{table}[H]
\centering
\caption{Сравнение наблюдаемых методов оптимизации}
\small
\begin{tabular}{lccc}
\toprule
Метод & GCC 13 & Clang 18 & TCC 0.9.27 \\
\midrule
Свёртка констант & + & + & + \\
Распространение констант & + & + & ограниченно \\
Распространение копий & + & + & ограниченно \\
Арифметические упрощения & + & + & частично \\
Удаление общих подвыражений & + & + & ограниченно \\
Снижение мощности & + & + & частично \\
Удаление лишних загрузок & + & + & ограниченно \\
Удаление недостижимого кода & + & + & слабее \\
Удаление лишних присваиваний & + & + & слабее \\
Использование регистров & + & + & + \\
Сжатие переходов & + & + & ограниченно \\
Удаление ненужных циклов & + & + & не выявлено \\
Разворачивание циклов & зависит от участка & зависит от участка & не выявлено \\
Слияние циклов & в тесте не произошло & в тесте не произошло & не произошло \\
\bottomrule
\end{tabular}
\end{table}

Отсутствие преобразования в конкретном тесте не означает, что компилятор вообще не поддерживает такую оптимизацию. Решение зависит от структуры программы, целевой архитектуры, уровня оптимизации и оценки стоимости преобразования.

% =========================================================
\section{Специфичные для сценария тесты}
% =========================================================

Для моделирования задач встраиваемого программного обеспечения были разработаны четыре дополнительных теста.

\subsection{Битовые маски}

\begin{lstlisting}[language=C,caption={Работа с битовыми регистрами}]
#include <stdio.h>
#include <stdint.h>

volatile uint32_t REG = 0;

void set_bits(uint32_t mask) {
    REG |= mask;
}

void clear_bits(uint32_t mask) {
    REG &= ~mask;
}

uint32_t test_bit(uint32_t mask) {
    return REG & mask;
}
\end{lstlisting}

Тест проверяет генерацию битовых операций и корректную работу с \texttt{volatile}.

\subsection{Кольцевой буфер}

\begin{lstlisting}[language=C,caption={Кольцевой буфер}]
#define SIZE 64

static uint8_t buf[SIZE];
static uint16_t head = 0, tail = 0;

void push(uint8_t b) {
    buf[head] = b;
    head = (head + 1) % SIZE;
}
\end{lstlisting}

Проверяется снижение мощности операции \texttt{\% 64}. Поскольку 64 является степенью двойки, компилятор может использовать битовую маску.

\subsection{CRC-32}

\begin{lstlisting}[language=C,caption={Основная операция CRC-32}]
crc = (crc >> 1) ^ (0xEDB88320 & -(crc & 1));
\end{lstlisting}

Тест содержит вложенные циклы, сдвиги, операции XOR и работу с массивом.

\subsection{Фильтр низких частот}

\begin{lstlisting}[language=C,caption={Целочисленный фильтр низких частот}]
#define ALPHA 8

int32_t lpf(int32_t x) {
    static int32_t y = 0;
    y = y + ((x - y) >> ALPHA);
    return y;
}
\end{lstlisting}

Проверяются арифметические сдвиги, работа со \texttt{static} и оптимизация небольших функций.

% =========================================================
\section{Результаты специфичных тестов}
% =========================================================

\begin{table}[H]
\centering
\caption{Размер секции \texttt{.text}, байт}
\small
\begin{tabular}{lrrrrrrr}
\toprule
Тест & \multicolumn{3}{c}{GCC} & \multicolumn{3}{c}{Clang} & TCC \\
\cmidrule(lr){2-4}\cmidrule(lr){5-7}
 & -O0 & -O2 & -Os & -O0 & -O2 & -Os & --- \\
\midrule
bitmask & 668173 & 668085 & 668021 & 676449 & 676281 & 676217 & 1046 \\
ringbuf & 668205 & 668125 & 667997 & 676345 & 676313 & 676249 & 957 \\
crc32   & 668237 & 668405 & 667917 & 676433 & 677025 & 676321 & 1052 \\
lpf     & 668045 & 667981 & 667917 & 676241 & 676225 & 676161 & 841 \\
\bottomrule
\end{tabular}
\end{table}

Для всех четырёх GCC-тестов минимальный размер среди исследованных режимов получен при \texttt{-Os}. Для CRC у GCC и Clang режим \texttt{-O2} увеличил размер кода относительно \texttt{-O0}, что показывает возможный компромисс между скоростью и объёмом программы.

\begin{table}[H]
\centering
\caption{Время выполнения специфичных тестов, с}
\small
\begin{tabular}{lrrrrrrr}
\toprule
Тест & GCC -O0 & GCC -O2 & GCC -Os & Clang -O0 & Clang -O2 & Clang -Os & TCC \\
\midrule
bitmask & 0.00 & 0.00 & 0.00 & 0.00 & 0.00 & 0.00 & 0.00 \\
ringbuf & 0.00 & 0.00 & 0.00 & 0.00 & 0.00 & 0.00 & 0.00 \\
crc32 & 0.02 & 0.02 & 0.00 & 0.02 & 0.01 & 0.02 & 0.02 \\
lpf & 0.00 & 0.00 & 0.00 & 0.00 & 0.00 & 0.00 & 0.00 \\
\bottomrule
\end{tabular}
\end{table}

Большинство дополнительных тестов оказались слишком короткими для точного сравнения времени выбранным способом. Поэтому значения 0.00 с не используются как доказательство превосходства одного режима над другим.

% =========================================================
\section{Ограничения эксперимента}
% =========================================================

При интерпретации результатов необходимо учитывать следующие ограничения:
\begin{enumerate}
    \item \texttt{optbench.c} является историческим синтетическим тестом;
    \item исследование проводится на x86-64, тогда как выбранный сценарий предполагает возможное использование ARM Cortex-M и RISC-V;
    \item статические бинарные файлы GCC и Clang включают значительный объём библиотечного кода;
    \item точность используемого измерения времени недостаточна для очень коротких тестов;
    \item отсутствие конкретной оптимизации в одном листинге не означает отсутствия такой возможности у компилятора вообще.
\end{enumerate}

% =========================================================
\section{Выводы}
% =========================================================

В ходе учебной практики было проведено сравнение оптимизаций компиляторов GCC 13.3.0, Clang 18.1.3 и Tiny C Compiler 0.9.27.

Основным компилятором выбран GCC. Исследование проводилось в Linux Mint 22.2 на архитектуре x86-64.

В качестве основной программы использовался тест \texttt{optbench.c}, построенный по методике Ф.~Н.~Хислея. Были рассмотрены сборки без оптимизации, с оптимизацией по скорости и с оптимизацией по размеру.

Эксперимент показал, что GCC и Clang способны:
\begin{itemize}
    \item удалять мёртвый и недостижимый код;
    \item удалять лишние присваивания;
    \item распространять константы и копии;
    \item упрощать арифметические выражения;
    \item оптимизировать циклы;
    \item устранять общие подвыражения;
    \item снижать мощность операций;
    \item сокращать цепочки переходов;
    \item эффективнее использовать регистры.
\end{itemize}

Особенно наглядными оказались функции \texttt{dead\_code} и \texttt{unnecessary\_loop}, в которых GCC и Clang при \texttt{-O2} фактически устраняют вычисления, не влияющие на результат.

TCC использует существенно более простую модель генерации кода и не демонстрирует такого количества глобальных оптимизаций, хотя некоторые локальные преобразования выполняет.

Исследование размеров показало, что \texttt{-O2} не всегда уменьшает объём машинного кода. В ряде случаев оптимизация производительности увеличивает размер программы. Режим \texttt{-Os}, напротив, ориентирован на уменьшение размера и в большинстве исследованных GCC-тестов дал минимальное значение.

Дополнительно были разработаны тесты битовых масок, кольцевого буфера, CRC-32 и фильтра низких частот. Они отражают операции, типичные для системного и встраиваемого программирования.

Таким образом, качество оптимизации нельзя оценивать только по одному параметру. Необходимо одновременно учитывать размер машинного кода, набор выполняемых инструкций и фактическую производительность программы. Для выбранного сценария GCC и Clang являются подходящими вариантами, а выбор конкретного уровня оптимизации должен зависеть от требований проекта к скорости и объёму кода.

% =========================================================
\section*{Список использованных источников}
\addcontentsline{toc}{section}{Список использованных источников}
% =========================================================

\begin{enumerate}
    \item Хислей Ф.~Н. Генерация высококачественного кода для программ, написанных на Си. PC Tech Journal, 1988.
    \item GNU Project. GCC Online Documentation. Optimize Options.
    \item LLVM Project. Clang Compiler User's Manual.
    \item Bellard F. Tiny C Compiler Reference Documentation.
    \item GNU Binutils Documentation: \texttt{objdump}, \texttt{size}.
    \item Linux manual pages: \texttt{gcc}, \texttt{clang}, \texttt{time}.
    \item Исходный файл \texttt{optbench.c}, предоставленный для выполнения учебной практики.
\end{enumerate}

% =========================================================
\appendix
\section{Основные команды эксперимента}
% =========================================================

\begin{lstlisting}[language=bash]
# GCC
gcc -O0 -static -DNO_ZERO_DIVIDE=1 optbench.c -o gcc_O0
gcc -O2 -static -DNO_ZERO_DIVIDE=1 optbench.c -o gcc_O2
gcc -Os -static -DNO_ZERO_DIVIDE=1 optbench.c -o gcc_Os

# Clang
clang -O0 -static -DNO_ZERO_DIVIDE=1 optbench.c -o clang_O0
clang -O2 -static -DNO_ZERO_DIVIDE=1 optbench.c -o clang_O2
clang -Os -static -DNO_ZERO_DIVIDE=1 optbench.c -o clang_Os

# TCC
tcc -DNO_ZERO_DIVIDE=1 optbench.c -o tcc_O0
\end{lstlisting}

\section{Рекомендуемая структура репозитория}

\begin{lstlisting}
compiler-optimization-study/
|-- README.md
|-- LICENSE
|-- .gitignore
|-- src/
|   |-- optbench.c
|   |-- bitmask.c
|   |-- ringbuf.c
|   |-- crc32.c
|   `-- lpf.c
|-- asm/
|   |-- gcc/
|   |-- clang/
|   `-- tcc/
|-- results/
|   |-- sizes/
|   `-- timings/
`-- .github/
    `-- workflows/
        `-- ci.yml
\end{lstlisting}

\end{document}
